\documentclass[10pt,letterpaper]{article}
\usepackage{cornell-notes}

\title{Chapter 2: The Physical Layer}
\author{}
\date{}

\setDocTitle{Chapter 2: The Physical Layer}
\setDocSubtitle{Computer Networks}
\setDocVersion{v1.0}
\setDocStatus{Study Companion}
\setDocOwner{Computer Networks Cornell Notes}
\setDocAuthor{}
\setDocDate{\today}
\setDocSummary{Cornell-format study and review notes for Chapter 2: The Physical Layer.}

\setCornellCollection{Computer Networks Cornell Notes}
\setCornellUnitType{Chapter}
\setCornellUnitNumber{2}
\setCornellUnitTitle{The Physical Layer}
\setCornellCompanionLabel{Study and review companion}

\hypersetup{pdftitle={Chapter 2: The Physical Layer},pdfsubject={Cornell notes},pdfauthor={}}

\begin{document}
\maketitle
\notesection{Chapter overview}
\begin{overviewbox}
\textbf{Learning objectives}\begin{itemize}
  \item Use Fourier reasoning, Nyquist's limit, and Shannon's limit to connect bandwidth, noise, and data rate.
  \item Compare copper, fiber, radio, light, and satellite transmission media.
  \item Explain baseband and passband modulation plus frequency-, time-, and code-division multiplexing.
  \item Describe the structure of telephone, cellular, and cable access systems.
  \item Choose a physical-layer approach by relating reach, capacity, noise, sharing, and deployment cost.
\end{itemize}
\textbf{Prerequisite concepts}\quad Signals, logarithms, frequency, binary representation, and the layered model.\par
\textbf{Key themes}\quad Channel capacity; attenuation and distortion; spectrum; modulation; multiplexing; switching; last-mile access.\par
\textbf{Named mechanisms and standards}\quad Nyquist and Shannon limits, fiber optics, radio and microwave links, satellites, QAM, OFDM, FDM, TDM, CDMA, SONET/SDH, ADSL, cellular generations, and cable modems.
\end{overviewbox}

\notesection{Cornell cue-and-note record}

\begin{longtable}{@{}L{0.245\textwidth}|L{0.695\textwidth}@{}}
\rowcolor{CNavy}\color{white}\bfseries Cue / recall prompt &
\color{white}\bfseries Notes \\
\endfirsthead
\rowcolor{CNavy}\color{white}\bfseries Cue / recall prompt &
\color{white}\bfseries Notes (continued) \\
\endhead
\hline
\endfoot
\cellcolor{CCue}\textbf{2.1 Theoretical basis} & \begin{minipage}[t]{\linewidth}\raggedright Physical transmission is constrained by the frequency response of the medium and by noise.\par\begin{itemize}\item \textbf{2.1.1 Fourier analysis} A periodic signal can be represented as a sum of sinusoidal harmonics. A channel that suppresses high-frequency components changes the received waveform.\item \textbf{2.1.2 Bandwidth-limited signals} Finite bandwidth removes harmonics and causes distortion. Symbol rate, encoding, and channel bandwidth jointly determine whether a receiver can distinguish transmitted states.\item \textbf{2.1.3 Maximum channel data rate} For a noiseless channel of bandwidth $B$ Hz using $V$ signal levels, Nyquist gives $R_{\max}=2B\log_2 V$ bit/s. For a noisy channel, Shannon gives $C=B\log_2(1+S/N)$ bit/s, where $S/N$ is a power ratio. These are upper bounds, not promised operating rates.\end{itemize}\end{minipage} \\ \hline
\cellcolor{CCue}\textbf{2.2 Guided media} & \begin{minipage}[t]{\linewidth}\raggedright Guided media confine energy to a physical path; their cost, bandwidth, attenuation, interference, and installation constraints differ.\par\begin{itemize}\item \textbf{2.2.1 Magnetic media} Physically transporting recorded media can deliver enormous bulk capacity but has very high latency; bandwidth and delay must be evaluated separately.\item \textbf{2.2.2 Twisted pairs} Twisting reduces electromagnetic interference. Categories, distance, and signaling determine achievable rate; telephone loops and Ethernet commonly use this inexpensive medium.\item \textbf{2.2.3 Coaxial cable} A shielded conductor offers greater bandwidth and noise immunity than basic twisted pair and is used in cable distribution systems.\item \textbf{2.2.4 Power lines} Power wiring can carry data but was not designed for high-frequency communication, so attenuation, noise, and changing electrical loads complicate operation.\item \textbf{2.2.5 Fiber optics} Light is guided through a core by refraction and total internal reflection. Single-mode fiber supports long distances and high rates; multimode fiber is easier to couple but experiences modal dispersion.\end{itemize}\end{minipage} \\ \hline
\cellcolor{CCue}\textbf{2.3 Wireless transmission} & \begin{minipage}[t]{\linewidth}\raggedright Unguided media radiate through space. Propagation, directionality, interference, licensing, and mobility shape their use.\par\begin{itemize}\item \textbf{2.3.1 Electromagnetic spectrum} Frequency bands differ in propagation and regulation. Usable bandwidth is a scarce resource allocated by national and international authorities, with some unlicensed bands.\item \textbf{2.3.2 Radio transmission} Lower-frequency radio often propagates broadly and penetrates obstacles, enabling omnidirectional coverage but also interference and limited frequency reuse.\item \textbf{2.3.3 Microwave transmission} Higher-frequency microwave is directional and often line-of-sight. Antenna alignment and atmospheric effects matter; narrow beams support spatial reuse.\item \textbf{2.3.4 Infrared transmission} Infrared is useful for short-range, room-confined links and generally does not penetrate walls, reducing cross-room interference.\item \textbf{2.3.5 Light transmission} Visible-light and laser links can carry high-rate signals along a narrow line of sight but are sensitive to obstruction and environmental conditions.\end{itemize}\end{minipage} \\ \hline
\cellcolor{CCue}\textbf{2.4 Communication satellites} & \begin{minipage}[t]{\linewidth}\raggedright \begin{itemize}\item \textbf{2.4.1 Geostationary satellites} A satellite above the equator can appear fixed and cover a large region, but its great distance produces substantial propagation delay.\item \textbf{2.4.2 Medium-Earth orbit} MEO systems trade lower delay and smaller coverage per satellite for a constellation larger than GEO.\item \textbf{2.4.3 Low-Earth orbit} LEO satellites reduce path length and delay but move relative to the ground, requiring many satellites, tracking, and handoff.\item \textbf{2.4.4 Satellites versus fiber} Satellites provide broadcast reach and rapid coverage; fiber provides very high point-to-point capacity and low error rates. Geography and traffic pattern determine the better fit.\end{itemize}\end{minipage} \\ \hline
\cellcolor{CCue}\textbf{2.5 Modulation and multiplexing} & \begin{minipage}[t]{\linewidth}\raggedright Modulation maps information to a waveform; multiplexing combines multiple logical streams on one physical resource.\par\begin{itemize}\item \textbf{2.5.1 Baseband transmission} Digital line codes map bits to signal levels on the medium. Clock recovery, direct-current balance, bandwidth, and noise immunity motivate schemes such as Manchester coding and block coding.\item \textbf{2.5.2 Passband transmission} A carrier is varied in amplitude, frequency, or phase. Combining amplitude and phase yields quadrature amplitude modulation, increasing bits per symbol when signal quality permits.\item \textbf{2.5.3 Frequency-division multiplexing} Users receive disjoint frequency bands separated by guards. Orthogonal frequency-division multiplexing divides a channel into many closely spaced orthogonal subcarriers.\item \textbf{2.5.4 Time-division multiplexing} Users share the same frequency range in distinct time intervals. Synchronous systems reserve slots; statistical systems can assign capacity in response to demand.\item \textbf{2.5.5 Code-division multiplexing} Users transmit simultaneously with distinct chip sequences. Correlation with the intended code recovers that user's data when codes and power are managed appropriately.\end{itemize}\end{minipage} \\ \hline
\cellcolor{CCue}\textbf{2.6 Public switched telephone network} & \begin{minipage}[t]{\linewidth}\raggedright \begin{itemize}\item \textbf{2.6.1 Structure} Local loops connect subscribers to end offices; trunks connect switching offices. The hierarchy has evolved toward digital transmission and switching.\item \textbf{2.6.2 Politics of telephones} Regulation, monopoly history, competition, and standardization influenced technical architecture and service boundaries.\item \textbf{2.6.3 Local loop} Modems carry data over voice-grade loops; asymmetric digital subscriber line divides loop spectrum into bands and adapts subcarriers to line quality; fiber moves optical capacity closer to users.\item \textbf{2.6.4 Trunks and multiplexing} Digital trunks aggregate many channels with time-division hierarchies; SONET/SDH provides synchronous framing, multiplexing, and operational support for optical transport.\item \textbf{2.6.5 Switching} Circuit switching reserves an end-to-end path for a session, message switching stores complete messages, and packet switching stores and forwards smaller units over shared links.\end{itemize}\end{minipage} \\ \hline
\cellcolor{CCue}\textbf{2.7 Mobile telephone system} & \begin{minipage}[t]{\linewidth}\raggedright Cellular design gains capacity through small cells, frequency reuse, mobility management, and handoff.\par\begin{itemize}\item \textbf{2.7.1 First generation} Analog voice cellular systems divided areas into cells and reused frequencies, but offered limited capacity and weak protection.\item \textbf{2.7.2 Second generation} Digital voice, time or code sharing, SIM-based identity, and text messaging improved capacity and services.\item \textbf{2.7.3 Third generation} Wideband code-based radio and packet data increased data capability while retaining cellular mobility and operator control.\end{itemize}\end{minipage} \\ \hline
\cellcolor{CCue}\textbf{2.8 Cable television} & \begin{minipage}[t]{\linewidth}\raggedright \begin{itemize}\item \textbf{2.8.1 Community antenna television} Early cable systems distributed received broadcast signals over coaxial trees to improve reception.\item \textbf{2.8.2 Internet over cable} Upgraded hybrid fiber-coax plants became two-way access networks in which many homes share downstream and upstream capacity.\item \textbf{2.8.3 Spectrum allocation} Frequency ranges are assigned to downstream television/data and upstream return traffic; upstream conditions are especially noisy.\item \textbf{2.8.4 Cable modems} A cable modem and headend termination system coordinate modulation, ranging, and shared upstream access.\item \textbf{2.8.5 ADSL versus cable} ADSL gives each subscriber a dedicated local loop to the office but shares capacity deeper in the network; cable shares the neighborhood medium earlier. Actual performance depends on line quality, users, and provisioning.\end{itemize}\end{minipage} \\ \hline
\cellcolor{CCue}\textbf{2.9 Summary} & \begin{minipage}[t]{\linewidth}\raggedright The physical layer converts data into signals, carries them through imperfect media, and shares costly channels. Capacity bounds, propagation, encoding, multiplexing, and switching explain the tradeoffs among access and backbone technologies.\end{minipage} \\ \hline
\end{longtable}

\begin{legacybox}Analog modems, early cellular generations, traditional telephone hierarchies, and the presented ADSL/cable deployments are retained to explain design evolution. They are not recommendations for new access deployments.\end{legacybox}

\notesection{Terminology and notation}

\begin{longtable}{@{}L{0.245\textwidth}|L{0.695\textwidth}@{}}
\rowcolor{CNavy}\color{white}\bfseries Cue / recall prompt &
\color{white}\bfseries Notes \\
\endfirsthead
\rowcolor{CNavy}\color{white}\bfseries Cue / recall prompt &
\color{white}\bfseries Notes (continued) \\
\endhead
\hline
\endfoot
\cellcolor{CCue}\textbf{Bandwidth $B$} & \begin{minipage}[t]{\linewidth}\raggedright The frequency range passed by a channel, measured in hertz; not identical to achieved bit rate.\end{minipage} \\ \hline
\cellcolor{CCue}\textbf{Baud} & \begin{minipage}[t]{\linewidth}\raggedright Number of transmitted symbols per second. One symbol may encode more than one bit.\end{minipage} \\ \hline
\cellcolor{CCue}\textbf{Signal-to-noise ratio} & \begin{minipage}[t]{\linewidth}\raggedright Power ratio $S/N$; decibels use $10\log_{10}(S/N)$.\end{minipage} \\ \hline
\cellcolor{CCue}\textbf{Attenuation} & \begin{minipage}[t]{\linewidth}\raggedright Reduction in signal strength as energy travels through a medium.\end{minipage} \\ \hline
\cellcolor{CCue}\textbf{Dispersion} & \begin{minipage}[t]{\linewidth}\raggedright Pulse spreading caused by frequency or propagation-mode differences.\end{minipage} \\ \hline
\cellcolor{CCue}\textbf{Modulation} & \begin{minipage}[t]{\linewidth}\raggedright Mapping information onto a baseband signal or carrier waveform.\end{minipage} \\ \hline
\cellcolor{CCue}\textbf{Multiplexing} & \begin{minipage}[t]{\linewidth}\raggedright Combining multiple logical channels for carriage over one physical resource.\end{minipage} \\ \hline
\cellcolor{CCue}\textbf{QAM} & \begin{minipage}[t]{\linewidth}\raggedright Quadrature amplitude modulation; symbols select combinations of carrier amplitude and phase.\end{minipage} \\ \hline
\cellcolor{CCue}\textbf{OFDM} & \begin{minipage}[t]{\linewidth}\raggedright Orthogonal frequency-division multiplexing; data is distributed across many orthogonal subcarriers.\end{minipage} \\ \hline
\cellcolor{CCue}\textbf{Propagation delay} & \begin{minipage}[t]{\linewidth}\raggedright Distance divided by propagation speed; independent of packet length.\end{minipage} \\ \hline
\end{longtable}

\notesection{Comparison matrix}

\begin{longtable}{@{}L{0.313\textwidth}L{0.313\textwidth}L{0.313\textwidth}@{}}
\toprule
\textbf{Technique} & \textbf{Shared dimension} & \textbf{Key tradeoff} \\ \midrule
FDM & Separate frequency bands & Guard bands and fixed spectrum allocation \\
TDM & Separate time intervals & Synchronization and unused reserved slots \\
CDM & Distinct spreading codes & Code correlation, power control, and interference \\
Circuit switching & Reserved end-to-end resources & Predictable allocation but inefficient idle periods \\
Packet switching & Statistical sharing by packets & Efficient bursts but variable queueing delay \\
\bottomrule
\end{longtable}
\notesection{Chapter synthesis}

\begin{overviewbox}\textbf{Cause-and-effect chain.} A source produces bits; an encoder maps them to symbols; modulation produces a waveform; the medium attenuates, distorts, and adds noise; a receiver estimates symbols; multiplexing lets multiple sources share the resource; switching connects resources into paths.\par
\textbf{Common confusions.} Hertz is not bit/s; baud is not always bit/s; transmission delay is not propagation delay; and Shannon's capacity is a bound, not an implementation recipe.\par
\textbf{Derived insight.} Increasing the number of signal levels raises bits per symbol only while the receiver can still distinguish the levels under the channel's noise and distortion.\end{overviewbox}

\notesection{Retrieval practice}

\begin{enumerate}
  \item For $B=3$ kHz and $V=4$, calculate the Nyquist upper bound for a noiseless channel.
  \item Explain how Shannon's $S/N$ term limits the benefit of adding signal levels.
  \item Why can shipped magnetic media have high bandwidth but poor interactive performance?
  \item Compare twisted pair, coaxial cable, and fiber in bandwidth, interference, and deployment cost.
  \item Why are microwave links directional while lower-frequency radio is often omnidirectional?
  \item Compare GEO and LEO satellite systems in coverage, delay, and constellation complexity.
  \item Distinguish baseband encoding from passband modulation.
  \item How do FDM, TDM, and CDM allocate a shared channel?
  \item Trace a subscriber's signal through the major portions of a telephone network.
  \item Why does packet switching suit bursty traffic better than fixed circuit allocation?
  \item How does cellular frequency reuse increase system capacity?
  \item Why can cable access performance vary with neighborhood activity?
\end{enumerate}
\begin{CornellSummaryBox}{Rapid-review Cornell summary}Channel bandwidth, signal levels, and noise bound data rate. Media differ in reach, attenuation, interference, and cost. Modulation maps data to waveforms; multiplexing shares media; switching connects channels. Always separate hertz, baud, bit/s, transmission delay, and propagation delay.\end{CornellSummaryBox}
\end{document}
